\section*{Chapitre \ref{ch:mx:execution-beyond-equities} --- Exécution au-delà des actions}

\begin{solution}{exo:mx:execution-beyond-equities:1}
Contrat à terme~: $0.7\times100\times\sqrt{50\text{M}/200\text{Md}}=1.11$ points de base. Obligation~: $0.7\times40\times\sqrt{50/20}=44.3$ points de base.
\end{solution}

\begin{solution}{exo:mx:execution-beyond-equities:2}
Vente de 200 contrats de l'échéance la plus proche~: 150 à 4\,999.75 et 50 à 4\,999.50, soit une moyenne de 4\,999.6875~; achat de 200 contrats de l'échéance suivante~: 100 à 5\,010.25 et 100 à 5\,010.50, soit une moyenne de
5\,010.375. Le spread est vendu à $-10.6875$ contre un prix milieu de $-10.25$~: 0.4375 points par contrat.
\end{solution}

\begin{solution}{exo:mx:execution-beyond-equities:3}
49.27, 46.26 et 45.54 points de base.
\end{solution}

\begin{solution}{exo:mx:execution-beyond-equities:4}
Un prix implied-in se construit à partir des carnets outright~; c'est donc un prix obtenu par legging, et le carnet de spreads direct offrait déjà le meilleur cours. La
liquidité implicite fait des carnets de spreads et d'outright un seul pool de liquidité, et permet d'exécuter le roll de façon atomique~: aucun risque que le prix de la seconde jambe bouge après la première.
\end{solution}

\begin{solution}{exo:mx:execution-beyond-equities:5}
Les cotations en streaming avec last look sont plus serrées, parce que le fournisseur peut refuser les demandes qui lui feraient perdre de l'argent~; le client ne paie cela que sur les
demandes rejetées (0.007 points de base ici). Cela ne serait plus avantageux si les rejets étaient fréquents, si le délai de détention était long, ou si les informations
des demandes rejetées étaient exploitées contre le client.
\end{solution}

\begin{solution}{exo:mx:execution-beyond-equities:6}
Le coût marginal du pool de liquidité est sa commission plus environ deux fois le mouvement de prix par dollar ($2x/R$)~; il égale le coût d'impact marginal de la bourse après environ
1.4~millions de dollars d'un pool de 1~milliard de dollars.
\end{solution}

\begin{solution}{exo:mx:execution-beyond-equities:7}
Deux dealers, 47.76 points de base~: une \omterm{def:mx:dark-pools-and-blocks:leakage}{fuite d'information} plus élevée fait que chaque dealer supplémentaire coûte davantage que la concurrence qu'il apporte.
\end{solution}

\begin{solution}{exo:mx:execution-beyond-equities:8}
Le prix payé est seulement celui de la demande acceptée~; les mouvements de prix perdus sur les demandes rejetées (et les nouvelles demandes au flux suivant) constituent le coût du last
look, soit 0.007 points de base en moyenne ici~: 0.087 au total, et non 0.080.
\end{solution}

\begin{solution}{pb:mx:execution-beyond-equities:1}
\textbf{1.} 50~millions de dollars représentent 0.025~\% de son volume quotidien.
\textbf{2.} Implied-in~: un prix de spread à partir des deux outrights~; implied-out~: un prix d'outright à partir du spread et de l'autre outright.
\textbf{3.} 0.88 points de base par legging, 0.45 par le carnet de spreads.
\textbf{4.} Exécution atomique et un seul pool de liquidité~; pas de meilleur prix ici.
\textbf{5.} Voir les définitions.
\textbf{6.} Transparence et divulgation~; le last look comme contrôle de risque, et non pour collecter de l'information.
\textbf{7.} Le fournisseur observe le mouvement pendant son délai de détention~; au-delà de son seuil, il rejette, et le client paie ce mouvement au flux suivant~: 0.007
points de base en moyenne.
\textbf{8.} Rejets de 7.7~\% des demandes~; 0.007 points de base.
\textbf{9.} Voir la définition.
\textbf{10.} 49.27 de marge, $-6.77$ de concurrence, 3.0 de \omterm{def:mx:dark-pools-and-blocks:leakage}{fuite d'information}~: 45.50.
\textbf{11.} Un quatrième dealer ajoute 1.46 points de base de concurrence, moins que ses 1.5 de \omterm{def:mx:dark-pools-and-blocks:leakage}{fuite d'information}.
\textbf{12.} 1.375~millions de dollars au pool et le reste à la bourse~; 0.36 points de base.
\textbf{13.} Sur les stratégies qui envoient beaucoup d'ordres et d'annulations (100 ordres en dix secondes sur une grande bourse), et non sur un calendrier en douze tranches.
\textbf{14.} \emph{Résultat nommé}~: coût total de 1.81 points de base sur le contrat à terme sur indice (le roll passant par le carnet de spreads, 25~\% de ce coût), 0.54 sur EUR/USD (last look,
1.3~\%), 45.50 sur l'obligation d'entreprise (la concurrence en RFQ, nette de la \omterm{def:mx:dark-pools-and-blocks:leakage}{fuite d'information}, fait économiser 8.3~\%) et 21.14 sur le bitcoin (la jambe passant par le pool fait économiser 1.7~\%).
\textbf{15.} La taille de l'ordre rapportée à la profondeur de chaque marché~: de 0.025~\% du volume d'une journée à 250~\%.
\textbf{16.} Un \omterm{def:mx:the-limit-order-book:orders}{ordre à cours limité} dans un carnet, une demande à un flux, une enchère entre dealers, un swap contre un pool.
\textbf{17.} Un ordonnanceur (\texttt{firm.acexec}) et un adaptateur par marché implémentant «~exécuter cette tranche~».
\textbf{18.} Les volumes et la marge des dealers obligataires.
\textbf{19.} Le trading électronique représentait 59~\% du volume des transactions de change en avril 2025, et un client pouvait choisir parmi plus de 15 plateformes multi-dealers publiées.
\textbf{20.} Ce qu'est un \omterm{def:mx:the-almgren-chriss-framework:parent}{ordre enfant}, et ce qu'il en coûte pour trouver la contrepartie.
\end{solution}

\begin{solution}{iq:mx:execution-beyond-equities:1}
En passant par le carnet de calendar spreads (et les prix implicites), et non par legging, sur les jours liquides de la période de roll, en surveillant la valeur du spread par rapport à son
carry théorique~; les risques sont le legging, une fenêtre de roll encombrée, et le spread qui évolue contre une grosse position.
\iqlookfor{Carnet de spreads~; timing~; risque de legging.}
\end{solution}

\begin{solution}{iq:mx:execution-beyond-equities:2}
Enregistrer, pour chaque demande, le mouvement de prix entre la demande et le rejet, ainsi que le prix obtenu lors de la nouvelle tentative~; le coût est la moyenne de ces mouvements sur l'ensemble des
demandes, auquel s'ajoute l'information que les demandes rejetées ont pu divulguer~; comparer avec le délai de détention et la symétrie publiés par le fournisseur.
\iqlookfor{Mesure des mouvements perdus~; divulgation.}
\end{solution}

\begin{solution}{iq:mx:execution-beyond-equities:3}
On ajoute des dealers jusqu'à ce que l'amélioration attendue apportée par un dealer supplémentaire (la concurrence) soit inférieure à la \omterm{def:mx:dark-pools-and-blocks:leakage}{fuite d'information} qu'il provoque~; cela dépend de la dispersion des prix des dealers,
de la \omterm{def:mx:dark-pools-and-blocks:leakage}{fuite d'information} par dealer perdant, et de la liquidité de l'obligation.
\iqlookfor{Le compromis~; ses paramètres.}
\end{solution}

\begin{solution}{iq:mx:execution-beyond-equities:4}
Un ordonnanceur demande une quantité à une échéance donnée~; l'adaptateur implémente «~exécuter cette tranche~» (placer et gérer des ordres, demander des flux, lancer une RFQ, faire un swap) et
publie les exécutions et les coûts dans un format commun. Conserver l'état propre à chaque marché à l'intérieur de l'adaptateur.
\iqlookfor{Séparation entre planification et mécanisme~; un enregistrement commun des exécutions.}
\end{solution}

\begin{solution}{iq:mx:execution-beyond-equities:5}
Égaliser les coûts marginaux~: l'impact marginal de la bourse face à la commission du pool augmentée du mouvement de prix ($2x/R$ pour un pool à produit constant), en incluant le gas et
le risque de front-running on-chain.
\iqlookfor{Coûts marginaux~; risques on-chain.}
\end{solution}

\begin{solution}{iq:mx:execution-beyond-equities:6}
Estimer le coût du débouclage (spread plus impact au volume de l'obligation, sur les jours nécessaires), le risque d'inventaire pendant sa détention, les couvertures
disponibles (contrats à terme sur taux, CDS sur indice), l'information du client et la concurrence sur la transaction~; coter un prix qui les couvre.
\iqlookfor{Coût du débouclage~; risque et couvertures~; concurrence.}
\end{solution}
